\section{Lowering Engine and Intermediate Representation}\label{sec:lowering-ir}

The \textit{Lowering Engine} is the most sophisticated component of the compiler.
It performs the transformation from the hierarchical, logically structured Annotated AST into a flat, chronologically
sorted list of events known as the Intermediate Representation (IR).
This process ``unrolls'' the recursive structure of the music and resolves all temporal logic into absolute beat
offsets, effectively stripping away all programming control flow.

\subsection{The Intermediate Representation (IR)}\label{detail-subsec:the-intermediate-representation-(ir)}

The output of the lowering engine is an \texttt{ir::Program} object.
This artifact represents the physical timeline of the composition.
Its primary building block is the \texttt{NoteEvent} structure:

\begin{lstlisting}[language=C++, caption={Structure of an atomic IR event.},label={detail-lst:lstlisting2}]
struct NoteEvent {
    int midi_note;          // Numeric pitch (0-127)
    double start_beat;      // Absolute position from track start
    double duration_beats;  // Temporal span in beats
    int velocity;           // MIDI intensity (default 100)
};
\end{lstlisting}

An important architectural detail is that the IR is entirely agnostic of the language's original control flow.
Whether a note was generated by a static \texttt{play} command, a \texttt{for} loop iteration, or a recursive
\texttt{pattern} call, it is represented in the IR as a simple, independent data point.

\subsection{Stateful Unrolling and Cursor Management}\label{detail-subsec:stateful-unrolling-and-cursor-management}

The unrolling engine is managed by the \texttt{LowererContext} class.
Unlike the semantic analyzer, which checks static types, the lowerer acts almost as an interpreter, executing the logic
and managing actual runtime values.
It maintains a ``cursor'' (a \texttt{double} passed by reference) tracking the current playhead position in beats.

The \texttt{lower\_statement} function utilizes \texttt{std::visit} to dispatch logic based on the statement kind:

\noindent \textbf{\texttt{PlayStatement}:} The lowerer evaluates the play target.
It determines the note duration (using the default, or overriding it if the \texttt{:} operator is present).
The start beat is set to the current cursor value, and the event is emitted.
Finally, the cursor is advanced by the duration.
If the \texttt{from} keyword is used, the start beat is overridden by the evaluated offset, and the cursor is
\textit{not} advanced.

\noindent \textbf{\texttt{LoopStatement} and \texttt{ForStatement}:} The engine evaluates the iteration count or
condition, then enters a loop in C++, recursively calling \texttt{lower\_block} for the body.
Because the \texttt{cursor} reference is passed down, the temporal state is naturally preserved across iterations,
placing notes end-to-end.

\noindent \textbf{\texttt{IfStatement}:} The lowerer evaluates the boolean condition into a concrete \texttt{true} or
\texttt{false}.
Only the matching branch is traversed, effectively performing ``compile-time pruning''.

\subsection{Dynamic Expression Evaluation}\label{detail-subsec:dynamic-expression-evaluation}

During the lowering phase, expressions are resolved into concrete values represented by the \texttt{ir::Value} variant.
These values can be scalars (\texttt{int}, \texttt{double}, \texttt{bool}) or musical structures (\texttt{NoteValue},
\texttt{RestValue}, \texttt{SequenceValue}, \texttt{ChordValue}).

The \texttt{LowererContext} maintains a \textit{Dynamic Scope Stack}: a vector of hash maps from identifier names to
\texttt{ir::Value} entries.
When a \texttt{VarDeclStatement} or assignment statement is encountered, the evaluated \texttt{ir::Value} is bound to
the identifier (types were already validated during semantic analysis).
This allows for dynamic algorithmic generation, where loop indices or accumulator variables can be used as multipliers
for note durations or offsets.

\subsection{Value Flattening and Complex Structures}\label{detail-subsec:value-flattening-and-complex-structures}

A significant challenge in the lowering engine is the conversion of complex musical types into atomic
\texttt{NoteEvent}s.
This is handled by the \texttt{ValueFlattener} module.

\subsubsection{Flattening Chords}
When a \texttt{ChordValue} is encountered, the flattener iterates over all its constituent notes and emits them using
the \textit{same} \texttt{start\_beat}.
The main lowering routine then advances the track's cursor by the \textit{maximum} duration found among the chord's
members, ensuring the next musical statement begins only after the entire chord has finished sounding.

\subsubsection{Pattern Temporal Reconstruction}
The most technically complex feature of the lowerer is how it preserves the ``temporal integrity'' of patterns.
When a \texttt{PatternCallExpression} is evaluated:

\noindent The lowerer snapshots the current cursor position: \texttt{start\_time = cursor}.

\noindent It pushes a new variable scope, binds the arguments to the parameter names, and unrolls the pattern's body.

\noindent It captures the final cursor state: \texttt{end\_time = cursor}.

\noindent It calculates the total temporal span as the difference between \texttt{end\_time} and \texttt{start\_time}.

To ensure that internal gaps and trailing rests within the pattern are respected, the lowerer rebuilds a
\texttt{SequenceValue} that intersperses the emitted notes with \texttt{RestValue} items matching the exact temporal
gaps.
This reconstructed value is then returned to the caller, allowing the pattern to be treated as a single, opaque rhythmic
block.

\subsection{Voice Parallelism and Cursor Isolation}\label{detail-subsec:voice-parallelism-and-cursor-isolation}

Parallel melodic lines are managed through the \texttt{voice} construct.
To implement parallelism without shared state conflicts or manual timing math, the lowerer uses
\textit{Cursor Isolation}:

\noindent Upon entering a \texttt{VoiceDefinition}, the lowerer captures the current outer cursor value.

\noindent It creates a local \textit{copy} called \texttt{voice\_cursor}.
If the voice specifies a \texttt{from} expression, the \texttt{voice\_cursor} is initialized to that value instead.

\noindent All statements unrolled within the voice modify only the \texttt{voice\_cursor}.

\noindent The generated \texttt{NoteEvent}s are appended to the track's master list, but upon exiting the
voice block, the
\texttt{outer\_cursor} is completely unmodified.

This isolation keeps the parent track and sibling voices aligned according to the compiler cursor model at
their starting beats, without requiring the author to calculate the intervening offsets manually.

\subsection{Implementation Layout and Control Flow}\label{detail-subsec:lowering-implementation-layout}

Lowering is implemented in \texttt{motivo\_lowering}.
The entry function \texttt{lowering/lower.cpp} destructures \texttt{analysis.program()} into header, globals, and
tracks, initializes an empty \texttt{ir::Program}, and constructs a \texttt{detail::LowererContext} bound to the
\texttt{AnalysisResult} and shared \texttt{DiagnosticsEngine}.

\noindent \textbf{Global initialization.} \texttt{lower\_header} copies tempo and time-signature fields into the IR
header.
\texttt{context.collect\_patterns(globals)} registers pattern symbols for global definitions; global
\texttt{VarDeclStatement} nodes are lowered inside a \texttt{LowererScopeGuard} so bindings exist before tracks run.

\noindent \textbf{Per-track lowering.} Each \texttt{TrackDefinition} is passed to \texttt{lower\_track\_definition}
(\texttt{detail/ast\_lowerers/lower\_definitions.cpp}), which collects track-local patterns, optionally validates
instrument constraints, and unrolls track items into an \texttt{ir::Track} with sorted \texttt{NoteEvent}s.

\noindent \textbf{Statement dispatch.} \texttt{detail/ast\_lowerers/lower\_statement.cpp} uses \texttt{std::visit} over
\texttt{StatementKind} to delegate to focused lowerers: \texttt{lower\_play\_statement.cpp},
\texttt{lower\_control\_flow\_statements.cpp} (\texttt{for}, \texttt{loop}, \texttt{if}), and
\texttt{lower\_declarative\_statements.cpp} (bindings and assignments).
\texttt{lower\_block} wraps each block in a \texttt{LowererScopeGuard} and concatenates emitted events.

\subsubsection{Expression Evaluation}

\texttt{detail/expression\_evaluator.cpp} dispatches on \texttt{ExpressionKind} via a C++20 \texttt{Literal} concept
for scalar and note literals.
Individual evaluators live under \texttt{detail/expression\_evaluators/}:
\texttt{evaluate\_identifier\_expression} resolves symbols through \texttt{LowererContext::lookup} using
\texttt{SymbolId} keys from semantic annotations;
\texttt{evaluate\_binary\_expression} and \texttt{evaluate\_unary\_expression} mirror runtime arithmetic;
\texttt{evaluate\_pattern\_call\_expression} snapshots the cursor, binds parameters in a fresh scope, executes the
pattern body through \texttt{execute\_block}, and reconstructs a \texttt{SequenceValue} preserving internal rests;
\texttt{evaluate\_sequence\_expression} and \texttt{evaluate\_chord\_expression} build composite \texttt{ir::Value}
structures.

\subsubsection{Value Flattening and Safety Limits}

\texttt{detail/value\_flattener.cpp} exposes \texttt{flatten\_value}, \texttt{flatten\_chord\_value}, and
\texttt{flatten\_sequence\_value}.
Chords emit parallel \texttt{NoteEvent}s at a shared \texttt{start\_beat}; sequences advance an internal cursor by each
item's effective duration (\texttt{get\_duration} handles notes, rests, and chords).

\texttt{LowererContext} enforces a hard cap of \texttt{MAX\_EVENTS = 20000} total emitted notes across the program via
\texttt{register\_events}; exceeding the limit reports a lowering error rather than producing an unusably large MIDI
file.
This guard complements the semantic phase, which validates structure but does not bound runtime unrolling depth.
